\chapter{Gasto federalizado}

\marginal{las participaciones suben, las aportaciones bajan}
Las dos grandes vías por las que el dinero federal llega a los estados se mueven
en direcciones contrarias. Las participaciones del Ramo 28 suben de 1,262,789.1 a
\textbf{1,340,210.9 millones de pesos}, un aumento real de 1.8\%. Las aportaciones
del Ramo 33 bajan de 984,484.5 a \textbf{979,951.8}, una caída real de 4.6\%. El
Ramo 25 cae 5.1\%.

La distinción importa porque las dos vías tienen naturaleza distinta: las
participaciones son ingreso propio de las entidades, de libre disposición; las
aportaciones están etiquetadas por ley a destinos concretos.

\textbf{El paquete 2025 transfiere más dinero libre y menos dinero etiquetado.}
Para las haciendas estatales es una mejora de margen; para los servicios que
financian los fondos etiquetados, salud y educación en primer lugar, es lo
contrario.

\begin{table}[htbp]
\centering
\caption{Gasto federalizado por vía, fondo y entidad federativa (mdp)}
\label{tab:C11_1}
\small
\begin{tabular}{lSSSSS}
\toprule
{Concepto} & {2024} & {2025} & {Dif. nominal} & {Dif. real} & {Var. real \%} \\
\midrule
Ramo 28 Participaciones & \num{1262789.1} & \num{1340210.9} & \num{77421.8} & \num{23121.9} & \num{1.8} \\
Ramo 33 Aportaciones & \num{984484.5} & \num{979951.8} & \num{-4532.7} & \num{-46865.5} & \num{-4.6} \\
Ramo 25 Previsiones y aportaciones & \num{82518.5} & \num{81668.0} & \num{-850.5} & \num{-4398.8} & \num{-5.1} \\
\multicolumn{6}{l}{\itshape Ramo 33 por fondo} \\
I013 & \num{454187.0} & \num{477160.5} & \num{22973.5} & \num{3443.4} & \num{0.7} \\
I005 & \num{116967.0} & \num{125352.3} & \num{8385.3} & \num{3355.7} & \num{2.8} \\
I004 & \num{101469.1} & \num{108743.4} & \num{7274.2} & \num{2911.1} & \num{2.8} \\
I002 & \num{135589.4} & \num{81220.5} & \num{-54369.0} & \num{-60199.3} & \num{-42.6} \\
I012 & \num{63908.9} & \num{68490.5} & \num{4581.6} & \num{1833.5} & \num{2.8} \\
I015 & \num{17901.9} & \num{18591.1} & \num{689.2} & \num{-80.6} & \num{-0.4} \\
I006 & \num{17092.9} & \num{18318.3} & \num{1225.4} & \num{490.4} & \num{2.8} \\
I003 & \num{13996.1} & \num{14999.4} & \num{1003.4} & \num{401.5} & \num{2.8} \\
I007 & \num{12842.0} & \num{13762.6} & \num{920.6} & \num{368.4} & \num{2.8} \\
I016 & \num{12880.5} & \num{13376.4} & \num{495.9} & \num{-58.0} & \num{-0.4} \\
I014 & \num{11823.4} & \num{12278.6} & \num{455.2} & \num{-53.2} & \num{-0.4} \\
I011 & \num{9210.9} & \num{9565.5} & \num{354.6} & \num{-41.4} & \num{-0.4} \\
I008 & \num{7223.6} & \num{7741.5} & \num{517.9} & \num{207.2} & \num{2.8} \\
I009 & \num{5798.0} & \num{6561.8} & \num{763.8} & \num{514.5} & \num{8.5} \\
I010 & \num{3593.8} & \num{3789.5} & \num{195.7} & \num{41.2} & \num{1.1} \\
\multicolumn{6}{l}{\itshape Ramos 28 y 33 por entidad federativa} \\
15 Estado de México & \num{273498.2} & \num{270082.7} & \num{-3415.6} & \num{-15176.0} & \num{-5.3} \\
09 Ciudad de México & \num{148678.8} & \num{151648.5} & \num{2969.8} & \num{-3423.4} & \num{-2.2} \\
30 Veracruz & \num{142372.8} & \num{144585.0} & \num{2212.2} & \num{-3909.8} & \num{-2.6} \\
14 Jalisco & \num{132403.9} & \num{138975.3} & \num{6571.4} & \num{878.0} & \num{0.6} \\
07 Chiapas & \num{109753.6} & \num{112439.1} & \num{2685.6} & \num{-2033.9} & \num{-1.8} \\
21 Puebla & \num{104449.0} & \num{105912.5} & \num{1463.5} & \num{-3027.8} & \num{-2.8} \\
19 Nuevo León & \num{96087.1} & \num{103125.0} & \num{7037.9} & \num{2906.2} & \num{2.9} \\
11 Guanajuato & \num{95121.1} & \num{100124.7} & \num{5003.6} & \num{913.4} & \num{0.9} \\
20 Oaxaca & \num{89463.3} & \num{90082.1} & \num{618.8} & \num{-3228.1} & \num{-3.5} \\
16 Michoacán & \num{83424.0} & \num{84729.2} & \num{1305.2} & \num{-2282.1} & \num{-2.6} \\
12 Guerrero & \num{77079.1} & \num{78709.2} & \num{1630.1} & \num{-1684.3} & \num{-2.1} \\
08 Chihuahua & \num{67535.0} & \num{71184.4} & \num{3649.3} & \num{745.3} & \num{1.1} \\
28 Tamaulipas & \num{65766.9} & \num{66753.7} & \num{986.7} & \num{-1841.2} & \num{-2.7} \\
02 Baja California & \num{63556.9} & \num{66138.8} & \num{2581.9} & \num{-151.0} & \num{-0.2} \\
13 Hidalgo & \num{58105.9} & \num{58771.3} & \num{665.4} & \num{-1833.2} & \num{-3.0} \\
25 Sinaloa & \num{53817.2} & \num{55650.9} & \num{1833.7} & \num{-480.5} & \num{-0.9} \\
26 Sonora & \num{52912.9} & \num{55003.0} & \num{2090.1} & \num{-185.2} & \num{-0.3} \\
05 Coahuila & \num{52093.1} & \num{54636.2} & \num{2543.1} & \num{303.1} & \num{0.6} \\
24 San Luis Potosí & \num{51145.5} & \num{52114.7} & \num{969.1} & \num{-1230.1} & \num{-2.3} \\
27 Tabasco & \num{52382.0} & \num{51897.9} & \num{-484.1} & \num{-2736.5} & \num{-5.0} \\
34 No Distribuible Geográficamente & \num{29777.5} & \num{47598.6} & \num{17821.1} & \num{16540.7} & \num{53.3} \\
22 Querétaro & \num{41254.8} & \num{43718.8} & \num{2464.0} & \num{690.0} & \num{1.6} \\
31 Yucatán & \num{39715.7} & \num{41626.3} & \num{1910.5} & \num{202.8} & \num{0.5} \\
10 Durango & \num{34167.7} & \num{36744.5} & \num{2576.9} & \num{1107.6} & \num{3.1} \\
23 Quintana Roo & \num{33277.4} & \num{35021.8} & \num{1744.4} & \num{313.4} & \num{0.9} \\
17 Morelos & \num{33456.1} & \num{33741.3} & \num{285.2} & \num{-1153.4} & \num{-3.3} \\
32 Zacatecas & \num{32766.7} & \num{32602.4} & \num{-164.3} & \num{-1573.3} & \num{-4.6} \\
01 Aguascalientes & \num{26969.6} & \num{28666.1} & \num{1696.5} & \num{536.8} & \num{1.9} \\
\bottomrule
\end{tabular}
\par\vspace{2pt}\footnotesize Fuente: elaboración propia del ITED con los analíticos del PEF aprobado 2024 y 2025. Las diferencias en millones de pesos se reportan dos veces, nominal y real; la real está en pesos de 2025 con el deflactor de 4.3\% que declara el CGPE.
\end{table}


\section{Dentro del Ramo 33}

El fondo de nómina educativa, que es el mayor con diferencia, sube 0.7\% real. El
fondo de aportaciones para los servicios de salud, como mostró el capítulo 6, cae
42.6\%. Ésa es la asimetría interna del ramo: \textbf{la nómina educativa se
protege y el fondo de salud absorbe el ajuste.}

\section{El corte que este documento no publica}

El género acostumbra publicar gasto federalizado por habitante. \textbf{Este
documento no lo hace}, y la razón es que no tiene una fuente de población por
entidad federativa. Las proyecciones del Consejo Nacional de Población no están en
la carpeta de este ejercicio.

Se prefiere publicar el reparto por entidad en pesos, que es verificable, antes
que un per cápita cuyo denominador no podríamos declarar. Es la aplicación directa
de la quinta declaración: \textbf{toda cifra per cápita declara su denominador o
no se publica.}

\clearpage
